\documentstyle[% 


righttag% 


,11pt% 


,amssymb% 


,russian%               remove in Eng. 


,izvru%                 Set izven% in Eng. 


% ,izvru-ks             for brief communications 


]{amsart} 


 


%


 


\newcommand{\tts}[1]{} 


 


\theoremstyle{definition} 


\theorembodyfont{\rm} 


\newtheorem{defnnonum}{\hskip\parindent Определение} 


\newtheorem{notenonum}{\hskip\parindent Замечание} 


\renewcommand{\thedefnnonum}{} 


\renewcommand{\thenotenonum}{} 


%\hoffset=-15mm%                  For Eng. 


 


\setcounter{page}{67} 


\god{2000} 


\nomer{7~(458)} %                        Remove in Eng. 


%\nomer{Vol.\,44, No.\,7}%               Set in Eng. 


%\str{\pageref{begin}--\pageref{end}}%  Set in Eng. 


\udk{517.53} 


 


\author{\it Р.Ф.\,Шамоян} 


% NB:   ^^ remove in Eng. 


\title{Непрерывные функционалы и мультипликаторы степенных рядов одного 


класса аналитических в поликруге функций} 


 


\date{} 


% \pagestyle{myheadings}%         Set in Eng. 


\pagestyle{plain} %              Remove in Eng. 


\begin{document} 


% \thispagestyle{plain}%          Set in Eng. 


\maketitle 


\label{begin} 


 


\section*{Введение} 


 


Пусть $U^n=\{z=(z_1,\dotsc,z_n),\ |z_j|<1,\ j=1,\dotsc,n\}$ --- 


единичный полидиск $n$-мерного комплексного пространства $\Bbb C^n$, 


$\Bbb T^n=\{z:|z_j|=1,\ j=1,\dotsc,n\}$ --- его остов, $H(U^n)$ --- 


множество всех голоморфных в $U^n$ функций. В дальнейшем будем 


использовать стандартные обозначения $I^n=[0,1]^n$, 


$(1-R)^\gamma=(1-R_1)^\gamma\dotsb(1-R_n)^\gamma$, 


$-\infty<\gamma<\infty$, $R_j\in(0,1)$, $j=1,\dotsc,n$, $s 


z=(s_1z_1,\dotsc,s_n z_n)$, $s\in\Bbb R^n$, $z\in\Bbb C^n$, $d R=d 


R_1\dotsb d R_n$, $M_q(f,r)=\Big(\int\limits_{\Bbb T^n}|f(r\xi)|^q d 


m_n(\xi)\Big)^{1/q}$, $0<q\le\infty$, $d m_{2n}(z)$ --- мера Лебега на 


$U^n$, $H^s(U^n)$ --- класс Харди в поликруге $U^n$, $0<s\le\infty$. 


Для произвольных $0<p,q<\infty$, $-1<\alpha<\infty$ введем пространства 


$$ 


F^{p,q}_\alpha(U^n)=\bigg\{f\in H(U^n): 


\|f\|^p_{F^{p,q}_\alpha(U^n)}=\int_{\Bbb T^n}\bigg( 


\int_{I^n}|f(R\xi|^q(1-R)^\alpha d R\bigg)^{p/q} 


d m_n(\xi)<\infty\bigg\}, 


$$ 


где $m_n(\xi)$ --- нормированная мера Лебега на $\Bbb T^n$. 


 


Нетрудно показать, что при $\max(p,q)\le1$ пространство 


$F^{p,q}_\alpha$ является полным метрическим пространством с метрикой 


$\rho(f,g)=\|f-g\|^{\min\{p,q\}}$, а при $\min\{p,q\}\ge1$ относительно 


нормы $\|f\|_{F^{p,q}_\alpha}$ пространство $F^{p,q}_\alpha$ банахово. 


Заметим, что $F^{p,p}_\alpha=A^p_\alpha$, где $A^p_\alpha(U^n)$ --- 


известные классы Бергмана--Джрбашяна \cite{1}, а 


$F^{p,2}_\alpha(U)=H^p_\alpha$, $0<p<\infty$ \cite{2}, где $H^p_\alpha$ 


--- классы Харди--Соболева. Некоторые свойства введенных пространств и 


их гладкостных аналогов при $1<p,q<\infty$ исследовались в \cite{3} и 


\cite{4}. 


 


Цель статьи --- изучение некоторых свойств пространств $F^{p,q}_\alpha$ 


при $0<p,q\le1$, $-1<\alpha<\infty$. 


 


Пусть $0<p,q\le1$, $-1<\alpha<\infty$. Обозначим через $S^{p,q}_\alpha$ 


множество всех голоморфных в $U^n$ функций таких, что для любого 


$\beta>\frac{\alpha+1}{q}+\frac{1}{p}$ 


$$ 


\|f\|_{S^{p,q}_\alpha}=\sup_{z\in U^n}\big\{|D^{\beta+1}g(z)| 


(1-|z|)^{\beta+2-\frac{(\alpha+1)}{q}-\frac{1}{p}}\big\}<\infty, 


$$ 


где 


$$ 


D^\beta g(z)=\sum_{|k|\ge0} 


\frac{\Gamma(k+\beta+1)}{\Gamma(\beta+1)\Gamma(k+1)}c_k z^k,\ \;


\Gamma(k+\beta+1)=\prod^n_{j=1}\Gamma(k_j+\beta+1),\ \;


g(z)=\sum_{|k|\ge0}c_k z^k,\ \; \beta>-1, 


$$ 


--- дробная производная функции $g(z)$. Легко видеть, что относительно 


нормы $\|\cdot\|_{S^{p,q}_\alpha}$ пространство $S^{p,q}_\alpha$ 


является банаховым. 


 


\begin{defnnonum} 


Пусть $X$ и $Y$ --- подпространства $H(U^n)$. Скажем, что 


последовательность комплексных чисел $\{(c_k)_{k\in z^n_+}\}$ является 


мультипликатором из $X$ в $Y$, если для любой функции 


$f(z)=\sum\limits_{|k|\ge0}a_k z^k$, $f\in H(U^n)$, функция 


$g(z)=\sum\limits_{|k|\ge0}c_k a_k z^k$ принадлежит $Y$. 


\end{defnnonum} 


 


Множество таких последовательностей обозначим через $M_T(X,Y)$. 


Отметим, что задача описания мультипликаторов для различных пар 


пространств $X$ и $Y$ решалась в работах многих авторов (см. 


\cite{5}--\cite{7}). 


 


\begin{thm} 


$1)$ Пусть $\Phi$ --- линейный непрерывный функционал на 


$F^{p,q}_\alpha$, $0<p,q\le1$, $g(z)=\Phi(l_z)$, $l_z(w)=\frac{1}{1-z 


w}=\prod\limits^n_{k=1}\frac{1}{1-z_k w_k}$, $z\in U^n$. Тогда $g(z)\in 


S^{p,q}_\alpha$, функционал $\Phi$ представ\'им в виде 


\begin{equation} 


\Phi(f)=\lim_{\rho\to1-0}\frac{1}{(2\pi)^n} 


\int_{\Bbb T^n}f(\rho\xi)g(\rho\ov\xi)d m_n(\xi). 


\end{equation} 


 


$2)$ Обратно, любая функция $g\in S^{p,q}_\alpha$ по формуле $(1)$ 


порождает линейный непрерывный функционал на $F^{p,q}_\alpha$. Более 


того, справедливы оценки 


$$c_1(\alpha,p,q)\|\Phi\|\le\|g\|_{S^{p,q}_\alpha}\le 


c_2(\alpha,p,q)\|\Phi\|. 


$$


\end{thm} 


 


\begin{thm} 


Если $g\in H(U^n)$, $g(z)=\sum\limits_{|k|\ge0}c_k z^k$, то 


$$ 


\{(c_k)_{k\in z^n_+}\}\in M_T(F^{p,q}_\alpha(U^n),\,H^s(U^n)),\quad 


0<\max(p,q)\le s\le\infty,\quad 0<p,q\le1, 


$$ 


тогда и только тогда, когда 


\begin{equation} 


\sup_{r\in I^n}M_s(D^m g,r) 


(1-r)^{m+1-\frac{1}{p}-\frac{\alpha+1}{q}}<\infty 


\end{equation} 


для некоторого $m\in\Bbb N$, $m>\frac{1}{p}-1+\frac{\alpha+1}{q}$. 


\end{thm} 


 


В дальнейшем 


$$ 


A^{p,q}_\alpha(U^n)=\bigg\{f\in H(U^n):\|f\|^p_{A^{p,q}_\alpha}= 


\int_{I^n}M^p_q(f,r)(1-r)^\alpha d r<\infty\bigg\},\quad 


0<p,q<\infty.


$$ 


Пространства типа $A^{p,q}_\alpha$ изучались многими авторами (см. [5]). 


 


\begin{thm} 


Если $g\in H(U^n)$, $g(z)=\sum\limits_{|k|\ge0}c_k z^k$, то 


 


$1)$ $\{(c_k)_{k\in z^n_+}\}\in M_T(F^{p,q}_\alpha(U^n),\, 


F^{t,s}_\beta(U^n))$, $0<s\le t<\infty$, $0<p,q\le1$, $0<\max(p,q)\le 


s\le1$, $s\big(\frac{\alpha+1}{q}+\frac{1}{p}\big)<2$ тогда и только 


тогда, когда 


\begin{equation} 


\sup_{r\in I^n}M_t(f,r) 


(1-r)^{m+1-\frac{1}{p}+\frac{\beta+1}{s}-\frac{\alpha+1}{q}}<\infty 


\end{equation} 


для некоторого $m\in\Bbb N$, $m\ge\frac{2}{q}-1;$ 


 


$2)$ $\{(c_k)_{k\in z^n_+}\}\in M_T(A^{p,q}_\alpha(U^n),\, 


F^{t,s}_\beta(U^n))$, $0<p,q\le s\le1$, $0<s\le t<\infty$, 


$s\big(\frac{\alpha+1}{p}+\frac{1}{q}\big)<2$ тогда и только тогда, 


когда 


\begin{equation} 


\sup_{r\in I^n}M_t(f,r) 


(1-r)^{m+1-\frac{1}{q}+\frac{\beta+1}{s}-\frac{\alpha+1}{p}}<\infty 


\end{equation} 


для некоторого $m\in\Bbb N$, $m\ge\frac{2}{s}-1$. 


\end{thm} 


 


\begin{notenonum} 


Утверждение теоремы 1 при $p=q$ хорошо известно и установлено в 


\cite{1}. 


\end{notenonum} 


 


Приведем схемы доказательств теорем 1--3. Они основаны на следующих 


двух леммах. 


 


\begin{lem} 


Пусть $f,g\in H(U^n)$, $r\in I^n$, $\alpha>-1$. Тогда 


\begin{multline*} 


\frac{1}{(2\pi)^n}\int_{\Bbb T^n}f(r\ov t)g(r t)d m_n(t)=\bigg( 


\frac{\alpha}{\pi}\bigg)^n r^{-2\alpha}_1\dotsb r^{-2\alpha}_n \times\\


\times\int^{r_1}_0\dotsb\int^{r_n}_0\int_{\Bbb T^n}D^{\alpha+1} 


g(R\xi)f(R\ov\xi)\prod^n_{l=1}(r^2_l-R^2_l)^\alpha 


R\,d R\,d m_n(\xi). 


\end{multline*}  


\end{lem} 


 


\begin{lem} 


Пусть $0<\max(p,q)\le s\le1$, $-1<\alpha<\infty$, $f\in 


F^{p,q}_\alpha(U^n)$. Тогда имеет место неравенство 


$$ 


\bigg(\int_{U^n}|f(w)|^s(1-|w|)s\bigg(\frac{\alpha+1}{q}+ 


\frac{1}{p}\bigg)-2d m_{2n}(w)\bigg)^{1/s}\le 


c\|f\|_{F^{p,q}_\alpha(U^n)},


$$ 


где $c$ --- некоторая константа. 


\end{lem} 


 


Для доказательства необходимости условий (2)--(4) используется 


ограниченность оператора $T:(a_k)_{k\in z^n_+}\to(c_k a_k)_{k\in 


z^n_+}$, $f(z)=\sum\limits_{|k|\ge0}a_k z^k$, $f\in F^{p,q}_\alpha$, и 


лемма 2. При доказательстве достаточности условия (2) при $1<s<\infty$ 


используется лемма 1 и двойственность пространств $H^s(U^n)$ и 


$H^{s'}(U^n)$, $\frac{1}{s}+\frac{1}{s'}=1$. Доказательство 


достаточности условий (3) и (4) следует из лемм~1,~2, при этом 


используется оценка 


$$ 


\bigg(\int_{U^n}|f(w)|d m_{2n}(w)\bigg)^s\le c 


\int_{U^n}|f(w)|^s(1-|w|)^{2s-2}d m_{2n}(w),\quad 


f\in H(U^n),\quad 0<s\le1, 


$$ 


вытекающая из результатов \cite{1}. 


 


В заключение отметим, что при доказательстве теорем 1--3 существенную 


роль играет разбиение поликруга $U^n$ на диадические области $U_{j k}$ 


(см. \cite{1}). 


 


 


\begin{thebibliography}{9} 


 


\bibitem{1} 


Djrbashian A., Shamoyan F.A. {\em Topics in the theory of $A^p_\alpha$ 


spaces} // Teubner Texte zur Math. -- 1988. -- V.\,105. -- P.\,1--199.


\bibitem{2} 


Fefferman C., Stein E. {\em $H^p$ spaces of several variables} // Acta 


Math. -- 1972. -- V.\,129. -- \No\,3--4. -- P.\,137--173. 


\bibitem{3} 


Гулиев В.С., Лизоркин П.И. {\em Классы голоморфных и гармонических 


функций в поликруге в связи с их граничными значениями} // Тр. матем. 


ин-та РАН. -- 1993. -- Т.\,204. -- С.\,137--159. 


\bibitem{4} 


Aleksandrov A.B., Peller V.V. {\em Hankel operators and similarity to a 


contraction} // Intern. Math. Research Notices. -- 1996. -- \No\,9. -- 


P.\,263--275. 


\bibitem{5} 


Шведенко С.В. {\em Классы Харди и связанные с ними пространства 


аналитических функций в единичном круге, поликруге и шаре} // Итоги 


науки и техн. ВИНИТИ. Сер. матем. анализ. -- 1985. -- Т.\,23. -- 


С.\,3--124. 


\bibitem{6} 


Тригуб Р.М. {\em Мультипликаторы класса $H^p(\Bbb D^m)$ при 


$p\in(0,1)$, аппроксимативные свойства методов суммирования степенных 


рядов} // Докл. РАН. -- 1994. -- Т.\,335. -- \No\,6. -- С.\,697--699. 


\bibitem{7} 


Nawrocki M. {\em Multipliers, linear functionals and the Frech\'et 


envelope of the Smirnov class $N_*(U^n)$} // Trans. Amer. Math. Soc. -- 


1990. -- V.\,322. -- \No\,2. -- P.\,493--506. 


 


\end{thebibliography} 


 


\vskip 2cm 


 


\post{\it Брянский государственный}{\it Поступила} 


\post{\it педагогический университет}{07.10.1998} 


 


\label{end} 


\end{document} 


 


